\chapter{Mise en œuvre}

\section{Introduction}
Ce chapitre présente la mise en œuvre concrète de la fonctionnalité 
de recueil de feedbacks et de réclamations développée dans le cadre 
du stage chez Bewize. Après avoir défini les besoins dans le chapitre 
de cadrage et établi la conception dans le chapitre précédent, il 
s'agit maintenant de décrire en détail l'implémentation technique 
de chaque composant de la solution.

Ce chapitre suit la logique de développement adoptée pendant le stage, 
en présentant successivement : les technologies utilisées, la mise en 
place de l'environnement de développement, le développement du formulaire 
WebView en React JS, le développement du back-end Spring Boot, la mise 
en place de la base de données PostgreSQL, et enfin le développement 
du dashboard monitor.

%==========================================================
\section{Technologies utilisées}
%==========================================================

\subsection{Introduction aux choix technologiques}
Le développement de la fonctionnalité de recueil de feedbacks et de 
réclamations chez Bewize repose sur un ensemble de technologies modernes, 
complémentaires et éprouvées. Le choix de ces technologies a été guidé 
par les contraintes de l'environnement existant chez Bewize, les bonnes 
pratiques du développement logiciel moderne, et les exigences de 
performance et de maintenabilité du projet.

\subsection{React JS}

\begin{figure}[H]
    \centering
    
\includegraphics[width=0.22\textwidth]{Figures/reactjs.png}
    \caption{Logo React JS}
    \label{fig:reactjs}
\end{figure}

La figure \ref{fig:reactjs} présente le logo de React JS, bibliothèque 
JavaScript open-source développée et maintenue par Meta. React JS est 
utilisée dans ce projet pour développer deux interfaces distinctes : 
la page WebView contenant le formulaire de saisie, et le dashboard 
monitor permettant à l'équipe Bewize de consulter et traiter les retours.

\textbf{Utilisation dans le projet :}
\begin{itemize}
    \item \textbf{La page WebView :} formulaire de saisie de feedback 
    ou réclamation, accessible depuis l'application mobile
    \item \textbf{Le dashboard monitor :} interface interne de 
    consultation et traitement des retours utilisateurs
\end{itemize}

\textbf{Avantages pour le projet :}
\begin{itemize}
    \item Développement rapide grâce aux composants réutilisables
    \item Interface réactive et fluide grâce au DOM virtuel
    \item Intégration simple avec les APIs REST Spring Boot
    \item Responsive design optimal pour l'affichage mobile via WebView
\end{itemize}

\subsection{React Native}

\begin{figure}[H]
    \centering
    
\includegraphics[width=0.22\textwidth]{Figures/react_native.png}
    \caption{Logo React Native}
    \label{fig:react_native}
\end{figure}

La figure \ref{fig:react_native} présente le logo de React Native, 
framework open-source développé par Meta permettant de créer des 
applications mobiles natives pour iOS et Android. Dans ce projet, 
React Native constitue l'environnement dans lequel la page WebView 
développée est intégrée par l'équipe Bewize.

\subsection{Spring Boot}

\begin{figure}[H]
    \centering
    
\includegraphics[width=0.22\textwidth]{Figures/spring_boot.png}
    \caption{Logo Spring Boot}
    \label{fig:spring_boot}
\end{figure}

La figure \ref{fig:spring_boot} présente le logo de Spring Boot, 
framework Java open-source qui simplifie le développement d'applications 
back-end et de microservices. Spring Boot est utilisé dans ce projet 
pour développer l'ensemble du back-end de la solution, notamment les 
APIs REST consommées par la WebView et le dashboard monitor.

\textbf{Utilisation dans le projet :}
\begin{itemize}
    \item Exposition des endpoints REST API pour la WebView et le dashboard
    \item Réception, validation et traitement des données de feedback
    \item Persistance des données dans PostgreSQL via Spring Data JPA
    \item Gestion de la logique métier des feedbacks et réclamations
\end{itemize}

\textbf{Avantages pour le projet :}
\begin{itemize}
    \item Déjà utilisé par l'équipe Bewize, facilitant l'intégration
    \item Création rapide d'APIs REST robustes et bien structurées
    \item Documentation automatique via Swagger
    \item Intégration native avec Spring Data JPA
\end{itemize}

\subsection{Node.js}

\begin{figure}[H]
    \centering
    
\includegraphics[width=0.22\textwidth]{Figures/nodejs.png}
    \caption{Logo Node.js}
    \label{fig:nodejs}
\end{figure}

La figure \ref{fig:nodejs} présente le logo de Node.js, environnement 
d'exécution JavaScript côté serveur basé sur le moteur V8 de Chrome. 
Dans ce projet, Node.js est utilisé comme environnement d'exécution 
pour React JS et pour la gestion des dépendances via npm.

\textbf{Utilisation dans le projet :}
\begin{itemize}
    \item Environnement d'exécution pour React JS
    \item Gestion des dépendances via npm
    \item Scripts de build et de déploiement front-end
    \item Serveur de développement local
\end{itemize}

\subsection{Docker}

\begin{figure}[H]
    \centering
    
\includegraphics[width=0.22\textwidth]{Figures/docker.png}
    \caption{Logo Docker}
    \label{fig:docker}
\end{figure}

La figure \ref{fig:docker} présente le logo de Docker, plateforme de 
conteneurisation permettant d'empaqueter une application et toutes ses 
dépendances dans des conteneurs isolés et portables. Docker garantit 
un comportement identique dans tous les environnements d'exécution.

\textbf{Utilisation dans le projet :}
\begin{itemize}
    \item Conteneur pour l'application back-end Spring Boot
    \item Conteneur pour la base de données PostgreSQL
    \item Conteneur pour la page WebView React JS
    \item \texttt{docker-compose.yml} pour orchestrer l'ensemble
\end{itemize}

\subsection{Git et GitHub}

\begin{figure}[H]
    \centering
    
\includegraphics[width=0.22\textwidth]{Figures/github.png}
    \caption{Logo Git et GitHub}
    \label{fig:github}
\end{figure}

La figure \ref{fig:github} présente le logo de Git et GitHub. Git est 
le système de contrôle de version distribué le plus utilisé au monde, 
tandis que GitHub est la plateforme web qui héberge les dépôts Git et 
offre des fonctionnalités collaboratives avancées.

\textbf{Utilisation dans le projet :}
\begin{itemize}
    \item Versionnement du code source complet
    \item Gestion des branches selon Gitflow Bewize
    \item Pull requests et code reviews avec le Tech Lead
    \item Intégration CI/CD pour les tests automatisés
\end{itemize}

\subsection{PostgreSQL}

\begin{figure}[H]
    \centering
    
\includegraphics[width=0.22\textwidth]{Figures/postgresql.png}
    \caption{Logo PostgreSQL}
    \label{fig:postgresql}
\end{figure}

La figure \ref{fig:postgresql} présente le logo de PostgreSQL, système 
de gestion de base de données relationnelle open-source reconnu pour 
sa robustesse et ses excellentes performances. PostgreSQL est utilisé 
dans ce projet pour stocker l'ensemble des données liées aux feedbacks 
et réclamations.

\textbf{Utilisation dans le projet :}
\begin{itemize}
    \item Stockage des feedbacks et réclamations
    \item Gestion des statuts de traitement
    \item Requêtes de filtrage pour le dashboard monitor
\end{itemize}

\subsection{Récapitulatif des technologies}

\begin{table}[H]
\centering
\caption{Récapitulatif des technologies utilisées dans le projet}
\label{tab:recap_technologies}
\renewcommand{\arraystretch}{1.5}
\begin{tabular}{|p{3cm}|p{3cm}|p{6cm}|}
\hline
\textbf{Technologie} & \textbf{Couche} & \textbf{Rôle dans le projet} \\
\hline
React JS & Front-end web & Page WebView et dashboard monitor \\
\hline
React Native & Mobile & Application mobile Bewize \\
\hline
Spring Boot & Back-end & APIs REST, logique métier, persistance \\
\hline
Node.js & Environnement & Exécution JS et gestion npm \\
\hline
Docker & DevOps & Conteneurisation et déploiement \\
\hline
Git / GitHub & Versioning & Contrôle de version et collaboration \\
\hline
PostgreSQL & Base de données & Stockage des feedbacks et réclamations \\
\hline
\end{tabular}
\end{table}

Le tableau \ref{tab:recap_technologies} récapitule les sept technologies 
utilisées dans ce projet, organisées par couche applicative. Cette stack 
technique forme un ensemble cohérent et complémentaire, où chaque 
technologie joue un rôle précis et justifié dans l'architecture globale 
de la solution développée chez Bewize.

%==========================================================
\section{Mise en place de l'environnement de développement}
%==========================================================

\subsection{Prérequis techniques}
Avant de commencer le développement, il a été nécessaire de mettre en 
place un environnement de développement complet et fonctionnel.

\begin{table}[H]
\centering
\caption{Environnement de développement — Outils et versions}
\label{tab:env_dev}
\renewcommand{\arraystretch}{1.5}
\begin{tabular}{|p{4cm}|p{3cm}|p{6cm}|}
\hline
\textbf{Outil} & \textbf{Version} & \textbf{Utilisation} \\
\hline
Node.js & 18.x LTS & Environnement d'exécution JavaScript \\
\hline
npm & 9.x & Gestionnaire de paquets JavaScript \\
\hline
Java JDK & 17 LTS & Environnement d'exécution Spring Boot \\
\hline
Maven & 3.9.x & Gestionnaire de dépendances Java \\
\hline
Docker Desktop & Dernière version & Conteneurisation des services \\
\hline
Git & 2.x & Contrôle de version \\
\hline
Visual Studio Code & Dernière version & Éditeur front-end React JS \\
\hline
IntelliJ IDEA & Dernière version & IDE back-end Spring Boot \\
\hline
Postman & Dernière version & Test des APIs REST \\
\hline
\end{tabular}
\end{table}

Le tableau \ref{tab:env_dev} présente les neuf outils installés et 
configurés sur le poste de développement. Ces outils couvrent l'ensemble 
des besoins du projet, depuis l'environnement d'exécution JavaScript 
et Java jusqu'aux outils de test des APIs REST.

\subsection{Configuration de Docker}
L'ensemble des services de l'application a été conteneurisé avec Docker 
afin de garantir la cohérence entre les environnements de développement 
et de production. Le fichier \texttt{docker-compose.yml} orchestre 
les trois conteneurs principaux : le back-end Spring Boot, la base 
de données PostgreSQL et la page WebView React JS.

\begin{figure}[H]
    \centering
    \includegraphics[width=0.88\textwidth]{Figures/docker_compose1.png}
    \includegraphics[width=0.88\textwidth]{Figures/docker_compose2.png}
    \caption{Configuration complète du fichier docker-compose.yml}
    \label{fig:docker_compose}
\end{figure}

La figure \ref{fig:docker_compose} illustre la configuration complète 
du fichier \texttt{docker-compose.yml} utilisée pour orchestrer les 
trois conteneurs du projet. Cette configuration définit les services, 
les ports exposés, les variables d'environnement et les dépendances 
entre conteneurs, garantissant un démarrage automatique et ordonné 
de l'ensemble des services.

\subsection{Initialisation du dépôt Git}
Dès le début du projet, un dépôt Git a été initialisé sur GitHub 
pour versionner l'ensemble du code source. La structure de branches 
adoptée suit la méthodologie Gitflow :

\begin{itemize}
    \item \textbf{main :} branche principale contenant le code stable
    \item \textbf{develop :} branche de développement
    \item \textbf{feature/webview-form :} formulaire WebView
    \item \textbf{feature/dashboard :} dashboard monitor
\end{itemize}

\begin{figure}[H]
    \centering
    \includegraphics[width=0.88\textwidth]{Figures/git_branches.png}
    \caption{Structure des branches Git selon la méthodologie Gitflow}
    \label{fig:git_branches}
\end{figure}

La figure \ref{fig:git_branches} illustre la structure des branches Git 
adoptée pour ce projet selon la méthodologie Gitflow. Les branches de 
fonctionnalité ont été créées à partir de la branche develop et 
fusionnées dans main lors des jalons de livraison.

%==========================================================
\section{Développement du formulaire WebView}
%==========================================================

\subsection{Création du projet React JS}
La première étape du développement a consisté à créer le projet React JS 
qui hébergera le formulaire WebView. La structure du projet est organisée 
comme suit :

\begin{itemize}
    \item \textbf{src/components/ :} composants React JS réutilisables
    \item \textbf{src/services/ :} services de communication avec l'API
    \item \textbf{src/styles/ :} fichiers CSS et styles globaux
    \item \textbf{src/utils/ :} fonctions utilitaires partagées
\end{itemize}

\subsection{Développement du composant formulaire}
Le composant principal du formulaire WebView comprend les champs suivants :

\begin{itemize}
    \item \textbf{Type de retour :} sélecteur entre "Feedback" et 
    "Réclamation"
    \item \textbf{Informations utilisateur :} prénom, nom, email, 
    numéro étudiant et statut
    \item \textbf{Détails de la réclamation :} type, objet, description 
    et pièce jointe
    \item \textbf{Bouton d'envoi :} bouton de soumission du formulaire
\end{itemize}

\begin{figure}[H]
    \centering
    \includegraphics[width=0.6\textwidth]{Figures/webview_form_1.png}
    \includegraphics[width=0.6\textwidth]{Figures/webview_form_2.png}
    \caption{Interface complète du formulaire WebView de recueil de feedbacks}
    \label{fig:webview_form}
\end{figure}

La figure \ref{fig:webview_form} présente l'interface complète du 
formulaire WebView développé en React JS. La partie supérieure affiche 
les informations personnelles de l'utilisateur, tandis que la partie 
inférieure présente les détails de la réclamation ainsi que le bouton 
d'envoi. L'interface est épurée, intuitive et adaptée à un affichage 
sur écran mobile.

\subsection{Validation des données du formulaire}
Une validation côté client a été implémentée pour vérifier :

\begin{itemize}
    \item La présence obligatoire de tous les champs requis
    \item La présence obligatoire du message avec une longueur minimale
    \item La longueur maximale du message
\end{itemize}

\begin{figure}[H]
    \centering
    \includegraphics[width=0.6\textwidth]{Figures/webview_validation.png}
    \caption{Affichage des messages de validation du formulaire WebView}
    \label{fig:webview_validation}
\end{figure}

La figure \ref{fig:webview_validation} illustre le comportement du 
formulaire en cas de soumission avec des champs invalides ou vides. 
Des messages d'erreur explicites sont affichés en rouge directement 
sous chaque champ concerné, guidant l'utilisateur vers une saisie 
correcte et complète avant l'envoi.

\subsection{Connexion à l'API Spring Boot}
Une fois le formulaire validé, les données sont envoyées au back-end 
Spring Boot via une requête HTTP POST.

\begin{figure}[H]
    \centering
    
\includegraphics[width=0.6\textwidth]{Figures/webview_confirmation.png}
    \caption{Message de confirmation après l'envoi réussi du formulaire}
    \label{fig:webview_confirmation}
\end{figure}

La figure \ref{fig:webview_confirmation} présente le message de 
confirmation affiché à l'utilisateur après l'envoi réussi de son 
feedback ou de sa réclamation. Ce message inclut un numéro de référence 
unique permettant de retrouver le retour dans le dashboard monitor, 
ainsi que la date et l'heure d'envoi.

%==========================================================
\section{Développement du back-end Spring Boot}
%==========================================================

\subsection{Création du projet Spring Boot}
Le projet Spring Boot a été créé avec les dépendances suivantes :

\begin{itemize}
    \item \textbf{Spring Web :} pour la création des APIs REST
    \item \textbf{Spring Data JPA :} pour la persistance des données
    \item \textbf{PostgreSQL Driver :} pour la connexion à PostgreSQL
    \item \textbf{Spring Validation :} pour la validation des données
    \item \textbf{Lombok :} pour la réduction du code boilerplate
    \item \textbf{SpringDoc OpenAPI :} pour la documentation Swagger
\end{itemize}

\subsection{Développement de l'entité Feedback}
L'entité Feedback est la classe centrale du modèle avec les attributs :

\begin{itemize}
    \item \textbf{id :} identifiant unique auto-généré
    \item \textbf{type :} FEEDBACK ou RECLAMATION
    \item \textbf{message :} contenu textuel du retour
    \item \textbf{dateEnvoi :} date et heure d'envoi
    \item \textbf{statut :} NOUVEAU, EN\_COURS ou TRAITE
    \item \textbf{utilisateurId :} identifiant de l'utilisateur
\end{itemize}

\subsection{Développement du Controller}
Le FeedbackController expose les endpoints REST API suivants :

\begin{table}[H]
\centering
\caption{Endpoints REST API implémentés dans le FeedbackController}
\label{tab:endpoints}
\renewcommand{\arraystretch}{1.8}
\begin{tabular}{|p{1.8cm}|p{5cm}|p{2cm}|p{4cm}|}
\hline
\textbf{Méthode} & \textbf{Endpoint} & \textbf{Code retour} & 
\textbf{Description} \\
\hline
POST & /api/feedbacks & 201 Created & Créer un nouveau retour \\
\hline
GET & /api/feedbacks & 200 OK & Liste de tous les retours \\
\hline
GET & /api/feedbacks/\{id\} & 200 OK & Détail d'un retour \\
\hline
PUT & /api/feedbacks/\{id\}/statut & 200 OK & Mettre à jour le statut \\
\hline
GET & /api/feedbacks/stats & 200 OK & Statistiques des retours \\
\hline
\end{tabular}
\end{table}

Le tableau \ref{tab:endpoints} présente les cinq endpoints REST API 
développés dans le FeedbackController. L'endpoint POST retourne un 
code HTTP 201 lors de la création réussie d'un retour, tandis que 
les endpoints GET et PUT retournent un code HTTP 200.

%==========================================================
\section{Configuration de la base de données PostgreSQL}
%==========================================================

\subsection{Création des tables}

\begin{table}[H]
\centering
\caption{Structure de la table feedback dans PostgreSQL}
\label{tab:table_feedback}
\renewcommand{\arraystretch}{1.5}
\begin{tabular}{|p{3.5cm}|p{3cm}|p{2.5cm}|p{4cm}|}
\hline
\textbf{Colonne} & \textbf{Type SQL} & \textbf{Contrainte} & 
\textbf{Description} \\
\hline
id & BIGSERIAL & PRIMARY KEY & Identifiant auto-incrémenté \\
\hline
type & VARCHAR(20) & NOT NULL & FEEDBACK ou RECLAMATION \\
\hline
message & TEXT & NOT NULL & Contenu du retour \\
\hline
date\_envoi & TIMESTAMP & NOT NULL & Date et heure d'envoi \\
\hline
statut & VARCHAR(20) & NOT NULL & NOUVEAU, EN\_COURS, TRAITE \\
\hline
utilisateur\_id & BIGINT & NOT NULL & Identifiant utilisateur \\
\hline
\end{tabular}
\end{table}

Le tableau \ref{tab:table_feedback} décrit la structure de la table 
\texttt{feedback} créée automatiquement par Spring Data JPA dans la 
base de données PostgreSQL. Cette table comporte six colonnes permettant 
de stocker toutes les informations nécessaires à la gestion des retours 
utilisateurs, depuis leur création jusqu'à leur traitement complet 
par l'équipe Bewize.

%==========================================================
\section{Développement du dashboard monitor}
%==========================================================

\subsection{Présentation du dashboard}

\begin{figure}[H]
    \centering
    \includegraphics[width=0.9\textwidth]{Figures/dashboard_home.png}
    \caption{Page d'accueil du dashboard monitor Bewize}
    \label{fig:dashboard_home}
\end{figure}

La figure \ref{fig:dashboard_home} présente la page d'accueil du 
dashboard monitor développé en React JS. Cette interface affiche 
quatre indicateurs clés en temps réel : le nombre total de retours 
reçus, le nombre de réclamations, le nombre de retours en attente 
de traitement et le taux de traitement global.

\subsection{Page liste des retours}

\begin{figure}[H]
    \centering
    \includegraphics[width=0.9\textwidth]{Figures/dashboard_list.png}
    \caption{Page liste des feedbacks et réclamations dans le dashboard}
    \label{fig:dashboard_list}
\end{figure}

La figure \ref{fig:dashboard_list} présente la page de liste complète 
des feedbacks et réclamations reçus. Chaque retour est affiché sous 
forme de ligne dans un tableau paginé, avec son identifiant, son type, 
son objet, l'utilisateur émetteur, la date d'envoi et le statut de 
traitement.

\subsection{Fonctionnalité de filtrage}

\begin{figure}[H]
    \centering
    \includegraphics[width=0.9\textwidth]{Figures/dashboard_filter.png}
    \caption{Fonctionnalité de filtrage des retours dans le dashboard}
    \label{fig:dashboard_filter}
\end{figure}

La figure \ref{fig:dashboard_filter} illustre la fonctionnalité de 
filtrage avancé du dashboard monitor. Des boutons de filtre rapide 
permettent à l'équipe de basculer entre tous les retours, uniquement 
les feedbacks ou uniquement les réclamations, ainsi que de filtrer 
par statut.

\subsection{Tableau de bord statistiques}

\begin{figure}[H]
    \centering
    \includegraphics[width=0.9\textwidth]{Figures/dashboard_stats.png}
    \caption{Tableau de bord statistiques du dashboard monitor}
    \label{fig:dashboard_stats}
\end{figure}

La figure \ref{fig:dashboard_stats} présente le tableau de bord 
statistiques du dashboard monitor. Cette vue synthétique affiche la 
répartition des retours par type sous forme de graphique en donut, 
la répartition par statut sous forme de barres horizontales, et 
l'évolution du nombre de retours par semaine.

%==========================================================
\section{Tests de l'application}
%==========================================================

\subsection{Tests de l'API avec Postman}

\begin{figure}[H]
    \centering
    \includegraphics[width=0.9\textwidth]{Figures/postman_test.png}
    \caption{Tests des APIs REST avec Postman}
    \label{fig:postman_test}
\end{figure}

La figure \ref{fig:postman_test} illustre l'environnement de test 
Postman utilisé pour valider les APIs REST du projet. Chaque endpoint 
a été testé avec plusieurs scénarios : le cas nominal avec des données 
valides, les cas limites et les cas d'erreur avec des données invalides 
ou manquantes.

\subsection{Tests fonctionnels bout en bout}

\begin{table}[H]
\centering
\caption{Résultats des tests fonctionnels bout en bout}
\label{tab:tests}
\renewcommand{\arraystretch}{1.8}
\begin{tabular}{|c|p{5cm}|p{3cm}|p{3cm}|}
\hline
\textbf{ID Test} & \textbf{Scénario testé} & \textbf{Résultat attendu} & 
\textbf{Résultat obtenu} \\
\hline
T01 & Envoi d'un feedback valide depuis la WebView & 
Confirmation affichée, données en base & Conforme \\
\hline
T02 & Envoi d'une réclamation valide depuis la WebView & 
Confirmation affichée, données en base & Conforme \\
\hline
T03 & Envoi avec champs vides & 
Messages de validation affichés & Conforme \\
\hline
T04 & Consultation des retours dans le dashboard & 
Liste affichée correctement & Conforme \\
\hline
T05 & Filtrage par type dans le dashboard & 
Résultats filtrés correctement & Conforme \\
\hline
T06 & Changement de statut d'une réclamation & 
Statut mis à jour en base et affiché & Conforme \\
\hline
T07 & Affichage des statistiques & 
Données calculées et affichées & Conforme \\
\hline
\end{tabular}
\end{table}

Le tableau \ref{tab:tests} présente les résultats des sept tests 
fonctionnels bout en bout réalisés. Ces tests couvrent l'ensemble 
du parcours utilisateur, depuis l'envoi d'un retour depuis la WebView 
jusqu'à sa consultation et son traitement dans le dashboard monitor. 
Tous les tests ont été réalisés avec succès.

%==========================================================
\section{Conclusion}
%==========================================================

Ce chapitre a présenté en détail la mise en œuvre complète de la 
fonctionnalité de recueil de feedbacks et de réclamations développée 
dans le cadre du stage chez Bewize. L'implémentation a couvert 
l'ensemble des composants de la solution : les technologies utilisées, 
le formulaire WebView React JS, le back-end Spring Boot avec ses APIs 
REST, la base de données PostgreSQL et le dashboard monitor.

Chaque composant a été développé en suivant les bonnes pratiques du 
développement logiciel moderne, avec une attention particulière portée 
à la qualité du code, à la validation des données et à la gestion des 
erreurs. Le chapitre suivant présentera les résultats des tests 
réalisés et une évaluation globale de la solution développée.
